\documentclass[a4,10pt]{amsart}
\usepackage{fancyhdr}
\author{Robert Fisher}
\title{An Introduction to \emph{1 -- 116}}
\begin{document}
\maketitle

\emph{1 -- 116} (working title) is a collection of paintings\footnote{though
not all pieces are ``paintings'' in the traditional sense} created in
accordance with a pre-defined set of rules. 

Each painting is based on a randomly chosen number from one to 116. This
range comes from previous work with the periodic table of the
elements.

The first forty-nine paintings were made between 24 March and 11
May 2006 inclusive, with the numbers chosen by the artist's family and 
friends. The remaining paintings were begun on 16 March 2008, the
numbers being randomly generated by a Korn shell script running on a Sun
v210 server and e-mailed to the artist each day.

\section{The Rules of the Game}

\newtheorem{rules}{Rule}

\begin{rules}One painting is made every day, and it must be started and
finished between getting up and going to bed. The day's number is not known
in advance.
\end{rules}

\begin{rules}Paintings \emph{begin} with a number. They are not necessarily
literal representations of that number: it merely acts as the first link in
a chain of thought. The final meaning of the painting is not necessarily
associated with the initial number.
\end{rules}

\begin{rules}Each painting's support is a 29 x 29cm masonite board.
\end{rules}

\section{Understanding the Work}

Each individual painting is intended to work on its own, with a unique
meaning and aesthetic appeal. The look of the final assembled piece is
unknown and unconsidered, and the piece as a whole should have its own
macroscopic meaning, which is expected to differ from person to person.

If the viewer wishes, he may try to decode each painting, working back from
the final image to the number which was the seed of the idea. If that is
your intent, keep in mind:

\begin{itemize}
\item mathematical or arithmetic curiosities
\item years, ages, and dates
\item references to other artists' work and to popular culture
\item physics and, to a lesser extent, chemistry
\item codes and ciphers
\item the tools, media and materials used to create the piece
\item visual and linguistic puns
\item geometry
\item sometimes it's impossible to work out the original number
\end{itemize}

Ultimately, whatever you think it is, is what it is.

\end{document}
